\documentclass[../main.tex]{subfiles}
\begin{document}
\onecolumn
\section{Formal environment and agent definitions}

\subsection{Environment graph}
Let the warehouse environment be a directed graph
\[
  G_{\mathcal{W}} = (V, E), \qquad E \subseteq V \times V ,
\]
where $V$ is a finite set of vertices (traversable cells, intersections, stations)
and $E$ is a finite set of directed edges representing traversable segments.
Each edge carries a positive traversal cost
\[
  c : E \to \mathbb{R}_{>0}.
\]
For $v \in V$ we write $\mathcal{N}_{G}(v) = \{ u \in V : (v,u) \in E \}$ for the
set of successors of $v$, and
$B_r(v) = \{ u \in V : d_G(v,u) \le r \}$ for the $r$-hop neighbourhood of $v$
under the hop distance $d_G$.

\paragraph{Relation to CRAMP.}
CRAMP formulates MAPF in a discrete grid world rather than on an explicit graph
\cite{cramp}. The grid world is the special case in which
$V \subseteq \mathbb{Z}^2$ is the set of free cells,
$E = \{ (v,u) \in V \times V : \lVert v-u \rVert_1 = 1 \}$ is the
4-neighbour relation, and $c \equiv 1$. All definitions below therefore reduce
to the CRAMP setting when $G_{\mathcal{W}}$ is a unit-cost 4-connected grid.

\subsection{Agents, states and actions}
Let
\[
  \mathcal{N} = \{1, \dots, n\}
\]
be the finite set of agents, following CRAMP's agent indexing $a_i$,
$i \in \mathcal{N}$ \cite{cramp}. Each agent $i$ has a start vertex
$s_i \in V$ and a goal vertex $g_i \in V$, and its state is
\[
  x^i_t = \bigl( \mathrm{pos}^i_t, \sigma^i_t \bigr) \in \mathcal{X}^i = V \times \Sigma^i ,
\]
where $\mathrm{pos}^i_t \in V$ is the occupied vertex at discrete time
$t \in \mathbb{N}$ and $\sigma^i_t \in \Sigma^i$ encodes internal variables.
The global state space is
\[
  \mathcal{S} = \prod_{i \in \mathcal{N}} \mathcal{X}^i ,
  \qquad s_t = (x^1_t, \dots, x^n_t) \in \mathcal{S},
\]
and $\mathrm{pos}^i : \mathcal{X}^i \to V$ is the position projection.
CRAMP corresponds to $\Sigma^i = \emptyset$, i.e.\ the state of an agent is
exactly its vertex.

The action space of agent $i$ is the local move set
\[
  \mathcal{A}^i(v) = \mathcal{N}_G(v) \cup \{v\}, \qquad v = \mathrm{pos}^i_t ,
\]
i.e.\ move along one outgoing edge or wait. The joint action space is
$\mathcal{A} = \mathcal{A}^1 \times \dots \times \mathcal{A}^n$.
On a 4-connected grid this yields exactly the five discrete CRAMP actions
$\{\mathrm{N}, \mathrm{E}, \mathrm{S}, \mathrm{W}, \mathrm{stay}\}$
\cite{cramp}.

\subsection{Feasibility and collision constraints}
A joint action $a_t = (a^1_t, \dots, a^n_t)$ is \emph{admissible} in state $s_t$ if
\begin{align}
  a^i_t &\in \mathcal{A}^i(\mathrm{pos}^i_t) && \text{(edge feasibility)} \\
  a^i_t &\neq a^j_t && \forall i \neq j \ \text{(vertex conflict)} \\
  \neg \bigl( a^i_t = \mathrm{pos}^j_t \wedge a^j_t = \mathrm{pos}^i_t \bigr)
        && \forall i \neq j \ \text{(edge / swap conflict)} .
\end{align}
Transitions are $\mathrm{pos}^i_{t+1} = a^i_t$ for admissible joint actions.
A plan for agent $i$ is a path
$P_i = (p_0, p_1, \dots, p_{T_i})$ with $p_0 = s_i$, $p_{T_i} = g_i$ and
$(p_k, p_{k+1}) \in E \cup \{(v,v) : v \in V\}$; a solution is a set
$\{P_i\}_{i \in \mathcal{N}}$ satisfying (2)--(3) at every timestep.
The objective follows CRAMP's metrics, the makespan
$\max_{i} T_i$ and the flowtime $\sum_i T_i$ \cite{cramp}.

\subsection{Partial observability}
Each agent observes only a bounded region of $G_{\mathcal{W}}$. Define the
\emph{observation subgraph} of agent $i$ at time $t$ as the induced subgraph
\[
  G^i_t = G_{\mathcal{W}}\bigl[ B_r(\mathrm{pos}^i_t) \bigr]
\]
on the $r$-hop neighbourhood of its current vertex, together with a vertex
feature map
\[
  \phi^i_t : B_r(\mathrm{pos}^i_t) \to \mathbb{R}^{d_v}
\]
assigning to each visible vertex features for obstacle status, occupancy by
other agents, goals of visible agents, and local congestion statistics
(queue length, recent edge flow, expected travel time). The observation of
agent $i$ is the pair
\[
  o^i_t = \bigl( G^i_t, \phi^i_t, g_i \bigr) \in \mathcal{O}^i ,
  \qquad
  \mathcal{O} = \mathcal{O}^1 \times \dots \times \mathcal{O}^n .
\]
This generalises CRAMP's observation, which is a
$10 \times 10 \times 4$ tensor over the agent's field of view --
obstacles, positions of visible neighbours, goals of those neighbours, and the
agent's own goal channel -- plus its own goal position \cite{cramp}:
on a grid with $r$ chosen so that $|B_r(v)| = 100$ and $d_v = 4$,
$(G^i_t, \phi^i_t)$ is exactly that tensor.

\subsection{Communication graph and GNN policy}
Following CRAMP, local communication is defined on a second, agent-level graph
that must be distinguished from $G_{\mathcal{W}}$. At each timestep let
\[
  G^{\mathrm{c}}_t = \bigl( \mathcal{N}, E^{\mathrm{c}}_t \bigr),
  \qquad
  E^{\mathrm{c}}_t = \bigl\{ (i,j) \in \mathcal{N}^2 :
      \mathrm{pos}^j_t \in B_r(\mathrm{pos}^i_t) \bigr\},
\]
so nodes are agents and edges encode mutual visibility (proximity), with
adjacency matrix $A_t \in \{0,1\}^{N \times N}$ computed dynamically from each
agent's local observation, where $N$ is the maximum number of visible agents
\cite{cramp}.

Each agent encodes its observation into a node embedding
\[
  h^i_t = f^{\mathrm{enc}}_\theta ( o^i_t ) \in \mathbb{R}^{d}, \qquad d = 512,
\]
yielding a node embedding matrix $H_t \in \mathbb{R}^{N \times d}$
\cite{cramp}. A graph neural network $f_\theta$ with $L$ message-passing layers
operates on $(H_t, A_t)$,
\[
  H^{(\ell+1)}_t = \psi\Bigl( \hat{A}_t H^{(\ell)}_t W^{(\ell)} \Bigr),
  \qquad \ell = 0, \dots, L-1, \quad H^{(0)}_t = H_t ,
\]
where $\hat{A}_t$ is the gated (attention-weighted) adjacency controlling how
much influence each node exerts on its neighbours, $W^{(\ell)}$ are learnable
weights and $\psi$ a nonlinearity \cite{cramp}. The resulting per-agent graph
signal is pooled and mapped to the policy heads, giving a stochastic
decentralised policy
\[
  \pi^i_\theta \bigl( a^i_t \mid o^i_t, G^{\mathrm{c}}_t \bigr),
\]
with parameters $\theta$ shared across agents, together with auxiliary
heads for the value $V_\theta(o^i_t)$, action validity and blocking prediction
\cite{cramp}.

\subsection{Learning objective}
Agent $i$ maximises the discounted return
\[
  R^i_t = \sum_{k \ge 0} \gamma^k r^i_{t+k}, \qquad \gamma \in [0,1),
\]
with advantage
\[
  A(o^i_t, a^i_t) = \sum_{k=0}^{K-1} \gamma^k r^i_{t+k}
      + \gamma^K V_\theta(o^i_{t+K}) - V_\theta(o^i_t)
\]
as in CRAMP's actor--critic formulation \cite{cramp}. The reward follows
CRAMP's crowd-aware structure
\[
  r^i_t = R_m \mathbb{1}[\text{move}] + R_s \mathbb{1}[\text{stay}]
        + R_c \mathbb{1}[\text{collision}]
        + R_e \mathbb{1}[\text{all goals reached}] + R_n ,
\]
with $R_m = -0.3$, $R_c = -2.0$, $R_e = +20.0$ and $R_n$ the crowd-aware term
rewarding movement into less densely occupied regions \cite{cramp}. In the
graph setting the crowd-aware term is defined on $G_{\mathcal{W}}$ as a function
of the occupancy of $\mathcal{N}_G(v)$ rather than of grid neighbours.
Policy classes are written $\Pi^{\mathrm{arch}}$, and a policy
$\pi = (\pi^1, \dots, \pi^n) \in \Pi^{\mathrm{arch}}$ maps observations to
admissible joint actions.

\section{Extension: storage, depots and pickup-and-delivery}
The formalism above is complete for MAPF on $G_{\mathcal{W}}$. The following
layer is an \emph{extension} used only for the coupled-policy study; the graph
and GNN definitions are unchanged.

\subsection{Vertex roles}
Partition the vertex set as
\[
  V = V^{S} \,\dot\cup\, V^{D} \,\dot\cup\, V^{B} \,\dot\cup\, V^{O},
\]
where $V^{S}$ are storage vertices, $V^{D}$ depot / pickup--dropoff stations,
$V^{B}$ base or charging vertices, and $V^{O}$ transit-only vertices.
When roles are ignored, $V^{S} = V^{D} = V^{B} = \emptyset$ and the setting
collapses to CRAMP's start--goal MAPF.

\subsection{Items, orders and tasks}
Let $I$ be the finite set of item types and $Q$ the finite set of item units,
with type map $\iota : Q \to I$. A storage placement at time $t$ is a map
\[
  p_t : Q \to V^{S},
\]
and the induced storage configuration is
$s_t : I \to 2^{V^{S}}$, $s_t(\kappa) = \{ p_t(q) : \iota(q) = \kappa \}$.
Orders are sampled from a demand distribution $Z$; each realised order
$z \in \mathcal{Z}$ is a finite multiset of item units, $z \subseteq Q$.
Each order induces one or more tasks
\[
  \tau = \bigl( q, v^{\mathrm{pick}}, v^{\mathrm{del}} \bigr)
  \in T \subseteq Q \times V^{S} \times V^{D},
  \qquad v^{\mathrm{pick}} = p_t(q),
\]
and tasks arrive as a stream $(\tau_t)_{t \in \mathbb{N}}$ driven by $Z$.
Under this extension the goal $g_i$ of agent $i$ is no longer static but is set
by its currently assigned task, so that the MAPF objective is replaced by a
throughput objective.

\subsection{Coupled and decoupled policies}
A policy pair $(s, \pi)$ is \emph{decoupled}, $(s,\pi) \in \mathcal{P}^{\mathrm{dec}}$,
if $s_t \equiv s_0$ or evolves by exogenous rules independent of realised flows
on $G_{\mathcal{W}}$, and $\pi$ is trained under a fixed assumed
placement and demand. CRAMP and comparable decentralised MAPF learners fall in
$\mathcal{P}^{\mathrm{dec}}$: placement is fixed and only $\pi$ is learned.

A pair is \emph{closed-loop coupled}, $(s,\pi) \in \mathcal{P}^{\mathrm{cpl}}$,
if there exists a feedback map $F$ with
\[
  s_{t+1} = F(s_t, H_t),
  \qquad
  H_t = \bigl( (s_k, a_k, o_k, \tau_k) \bigr)_{k \le t},
\]
so that observed edge flows, queue lengths and travel times on $G_{\mathcal{W}}$
alter placement, which in turn alters the task-induced goal distribution and
hence the trajectories generated by $\pi_\theta$.

\subsection{Problem statement}
Given a fixed environment graph $G_{\mathcal{W}} = (V,E)$ with costs $c$,
item and order data $(I, Q, Z)$, agent set $\mathcal{N}$, and policy classes
$\mathcal{P}^{\mathrm{dec}}, \mathcal{P}^{\mathrm{cpl}}$, the thesis designs and
trains an RL-GNN controller $\pi_\theta \in \mathcal{P}^{\mathrm{cpl}}$ and
compares its induced performance functionals -- average travel distance per
completed task, average edge congestion, throughput, and robustness under
scaling of $|\mathcal{N}|$ and $|V|$ -- against decoupled baselines.

\end{document}